%==============================================================================
\section{Real-Life Examples --- Ví dụ học máy trong đời sống}
%==============================================================================

\begin{dongco}[title={ML không xa lạ như ta nghĩ}]
Học máy đã biến đổi nhiều ngành nhờ tự động hoá, dự đoán kết quả và phát hiện pattern trong dữ liệu lớn.
Ví dụ thực tế như trợ lý ảo và chatbot (Google Assistant, Siri, Alexa), hệ gợi ý, Tesla autopilot, IBM Watson for Oncology, \ldots
\end{dongco}

\begin{ynghia}[title={ML trong đời sống hằng ngày}]
Nhiều người nghĩ ML gắn với robot tương lai và rất phức tạp, nhưng thực tế mỗi ngày chúng ta dùng ML (có thể không để ý) như Google Maps, email, Alexa, \ldots
\end{ynghia}

\subsection{Danh sách ví dụ}

\begin{phuongphap}[title={Các ví dụ tiêu biểu}]
\begin{itemize}
  \item Trợ lý ảo và Chatbot (Virtual Assistants and Chatbots)
  \item Phát hiện gian lận trong ngân hàng và tài chính (Fraud Detection in Banking and Finance)
  \item Chẩn đoán và điều trị y tế (Healthcare Diagnosis and Treatment)
  \item Xe tự hành (Autonomous Vehicles)
  \item Hệ gợi ý (Recommendation Systems)
  \item Quảng cáo nhắm target (Target Advertising)
  \item Nhận dạng ảnh (Image Recognition)
\end{itemize}
\end{phuongphap}

\subsection{Virtual Assistants and Chatbots}

\begin{dongco}[title={NLP và trải nghiệm hội thoại}]
Xử lý ngôn ngữ tự nhiên (NLP) là một mảng của ML tập trung hiểu và sinh ngôn ngữ người.
NLP được dùng trong trợ lý ảo và chatbot như Siri, Alexa, Google Assistant để tạo trải nghiệm cá nhân hoá và hội thoại.
Thuật toán ML phân tích pattern ngôn ngữ và phản hồi truy vấn theo cách tự nhiên và chính xác.
\end{dongco}

\begin{ynghia}[title={Virtual assistants}]
Trợ lý ảo tương tác qua lệnh giọng nói, có thể thay công việc trợ lý cá nhân như gọi điện, lên lịch hẹn, hoặc đọc email.
Các trợ lý phổ biến: Alexa, Apple Siri, Google Assistant.
\end{ynghia}

\begin{ynghia}[title={Chatbots}]
Chatbot là chương trình ML được thiết kế để trò chuyện với người dùng.
Ứng dụng nhằm thay thế một phần công việc chăm sóc khách hàng: cung cấp thông tin, trả lời FAQ và hỗ trợ cơ bản.
\end{ynghia}

\subsection{Fraud Detection in Banking and Finance}

\begin{dongco}[title={ML cho an toàn và bảo mật}]
ML không chỉ để “làm dễ” mà còn để tăng an toàn: ví dụ phát hiện gian lận.
Thuật toán được huấn luyện trên dataset có hành vi gian lận/không mong muốn để học pattern và phát hiện khi chúng xảy ra trong tương lai.
\end{dongco}

\begin{ynghia}[title={Ví dụ thực tế}]
Thuật toán có thể phân tích giao dịch và nhận ra pattern chỉ ra gian lận.
Ví dụ: công ty thẻ tín dụng phát hiện giao dịch nghi gian lận và báo cho khách hàng theo thời gian thực.
Ngân hàng dùng ML để phát hiện rửa tiền, hành vi bất thường trong tài khoản, và phân tích rủi ro tín dụng.
PayPal là một ví dụ dùng ML để cải thiện giao dịch hợp lệ trên nền tảng.
\end{ynghia}

\subsection{Healthcare Diagnosis and Treatment}

\begin{dongco}[title={ML hỗ trợ y tế cá nhân hoá}]
Ứng dụng ML trong y tế rất đa dạng: cá nhân hoá điều trị, theo dõi bệnh nhân, chẩn đoán từ ảnh y khoa.
Kết hợp ML và y học nhằm tăng hiệu quả và tính cá nhân hoá.
\end{dongco}

\begin{ynghia}[title={Ví dụ}]
Thuật toán ML có thể phân tích dữ liệu y tế như X-ray, MRI, dữ liệu gene để hỗ trợ chẩn đoán.
ML cũng có thể gợi ý điều trị hiệu quả dựa trên tiền sử và gene của bệnh nhân.
Ví dụ: IBM Watson for Oncology dùng ML để phân tích hồ sơ y tế và đề xuất điều trị ung thư cá nhân hoá.
\end{ynghia}

\subsection{Autonomous Vehicles}

\begin{dongco}[title={Xe tự hành}]
Xe tự hành dùng ML để thay thế một phần người lái.
Xe được thiết kế để đến đích, tránh vật cản và phản ứng điều kiện giao thông.
Thuật toán ML phân tích dữ liệu từ cảm biến/camera để nhận ra vật cản và quyết định cách phản ứng.
\end{dongco}

\begin{ynghia}[title={Ví dụ}]
Xe tự hành có thể giúp giảm tai nạn và tăng hiệu quả.
Một số công ty: Tesla, Waymo, Uber đang dùng ML để phát triển xe tự lái.
Tesla Vision dùng camera/cảm biến và neural net processing để hiểu môi trường; AutoPilot là hệ thống hỗ trợ lái nâng cao.
\end{ynghia}

\subsection{Recommendation Systems}

\begin{dongco}[title={Cá nhân hoá ở quy mô lớn}]
Các nền tảng e-commerce/giải trí như Amazon và Netflix dùng hệ gợi ý (thuật toán ML) để đề xuất cá nhân hoá dựa trên lịch sử duyệt/xem.
Điều này tăng hài lòng khách hàng và tăng doanh thu.
\end{dongco}

\begin{ynghia}[title={Các ví dụ hệ gợi ý}]
\begin{itemize}
  \item Netflix: dùng lịch sử xem, tìm kiếm, rating để gợi ý phim/chương trình.
  \item Amazon: dùng sản phẩm đã xem, mua, thêm vào giỏ để gợi ý.
  \item Spotify: gợi ý bài hát/playlist theo lịch sử nghe, tìm kiếm, liked songs.
  \item YouTube: gợi ý video theo lịch sử xem, tìm kiếm, liked videos và nhiều yếu tố khác.
  \item LinkedIn: gợi ý việc làm/kết nối theo hồ sơ, kỹ năng, vị trí, ngành, \ldots
\end{itemize}
\end{ynghia}

\subsection{Target Advertising}

\begin{dinhnghia}[title={Quảng cáo nhắm target}]
Quảng cáo nhắm target dùng ML để khai thác insight từ dữ liệu nhằm “may đo” quảng cáo theo sở thích, hành vi và nhân khẩu học của cá nhân/nhóm.
\end{dinhnghia}

\subsection{Image Recognition}

\begin{dinhnghia}[title={Nhận dạng ảnh}]
Nhận dạng ảnh là ứng dụng computer vision, thường cần nhiều tác vụ CV như image classification, object detection và image identification.
Ứng dụng: nhận diện khuôn mặt, tìm kiếm hình ảnh (visual search), chẩn đoán y khoa, nhận dạng người, \ldots
\end{dinhnghia}

\begin{luuy}[title={Kết luận}]
Ngoài các ví dụ trên, ML còn dùng trong quản lý năng lượng, phân tích mạng xã hội, predictive maintenance, \ldots
Học máy có tiềm năng “cách mạng” nhiều ngành và cải thiện cuộc sống.
\end{luuy}
